% =============================================================================
% Section 3 -- Data, Preprocessing, and Annotation
% =============================================================================

\providecommand{\candidate}[1]{}
\providecommand{\dnote}[1]{}
\providecommand{\budget}[1]{}
\providecommand{\src}[1]{}
\providecommand{\authorinput}[1]{\textbf{\textless author input: #1\textgreater}}

\section{Data and Methods}

\yash{this section is WIP}

\label{sec:data}

\candidate{S3-lede}
Labelled data carry a reporting standard in social computing: an account of how
the dataset was built~\cite{gebru2021datasheets}, a codebook whose paradigm is
declared rather than left implicit~\cite{rottger2022paradigms}, annotators
described rather than assumed interchangeable~\cite{sap2022annotators}, and
machine labels checked against human ones before they carry
weight~\cite{pangakis2023validation}. We report against that standard and add
two practices to it. First, we chose a labeller by measurement rather than by
reputation. Four families of system, namely lexicon and encoder baselines,
open-weight language models with and without fine-tuning, a deployed commercial
classifier, and frontier commercial models, were scored against human gold
standards before any was adopted, and Table~\ref{tab:benchmark} reports that
comparison including the systems that failed outright. Second, every machine
agreement figure appears beside the human figure for the same construct.
Reporting a human baseline alongside machine labels is a minority practice at
this venue, and a machine figure cannot be read without knowing the precision of
the ruler it was scored against.

% -----------------------------------------------------------------------------
\subsection{Corpus and site}
\label{sec:data-site}
% -----------------------------------------------------------------------------

\subsubsection{Reddit as a site for interstate conflict discourse}
\label{sec:data-site-reddit}

\candidate{S3.1.1}
Reddit suits this question for reasons specific to the analyses that follow.
Pseudonymity lowers the cost of posting a nationalist sentiment under one's own
handle, and that register is what name-bound platforms suppress. Subreddits give
a boundary we can observe, because norms on Reddit are enforced at community
level and differ sharply between communities on the same
site~\cite{chandrasekharan2018hiddenrules}. That boundary is what turns a
comment left in one community by a participant from another into a measurable
event~\cite{datta2019interconflict}. Threading preserves who replied to whom.
Votes give a reception measure the communities produce themselves, and Reddit's
\texttt{controversiality} flag, which marks comments drawing roughly balanced
approval and disapproval, is what turns that reception into something we can
model. The rule behind the flag is Reddit's own and is unpublished, so our
outcome variable is a platform artefact as much as a record of
contestation~\cite{crawford2016flag}.

\subsubsection{Communities, alignment, and coverage}
\label{sec:data-communities}

% \begin{figure*}[t]
%   \centering
%   \includegraphics[width=\linewidth]{figures/s3_f1_corpus_composition}
%   \caption{Corpus composition and coverage, over the annotated corpus of
%   309,204 comments (2 January 2019 to 25 January 2026). Panel A gives comments
%   per community on a logarithmic axis, coloured by the alignment group each
%   community was assigned to. The India-aligned group spans six communities, from
%   a general-interest subreddit to a defence forum to a satire community; the
%   Pakistan-aligned group is one community, \texttt{r/pakistan}. A national
%   contrast drawn across them is therefore also a contrast between many
%   communities and one, which is why \S\ref{sec:results} reports every
%   community-level quantity within its group as well as across groups. Panel B
%   gives monthly comment volume on the same axis with the six crisis windows
%   shaded and labelled with their counts. Volume differs by more than three
%   orders of magnitude between the largest window and the quietest month.}
%   \label{fig:composition}
% \end{figure*}

\candidate{S3.1.2}
We selected communities by running the collection keywords as web queries and
taking the subreddits whose posts ranked highest, then checked that ranking
against the platform itself. On the Pakistan side we compared subscriber counts
and keyword-matching post volume across \texttt{r/pakistani},
\texttt{r/karachi}, \texttt{r/Lahore}, \texttt{r/pkmigrate},
\texttt{r/PakistaniMilitary} and \texttt{r/PakistanDiscussions}. None cleared a
volume threshold worth adding, which is itself a finding about the setting
rather than only about our procedure: \texttt{r/pakistan} appears to be the one
active Pakistan-aligned community for this discourse. We ran that check on the
Pakistan side alone, because the asymmetry it addresses runs in one direction,
and we did not run the equivalent check on the India side. The India-aligned
group is therefore plural by observation and the Pakistan-aligned group singular
by observation, and neither is plural or singular by design.

\candidate{S3.1.3}
Figure~\ref{fig:composition}A gives the result. Fourteen communities were
collected: India-aligned holds 185,177 comments across six communities, neutral
66,732 across three, Pakistan-aligned 53,308 in one, and a residual four
communities hold 3,987 between them. Those four, \texttt{r/abcdesis},
\texttt{r/askreddit}, \texttt{r/kashmir} and \texttt{r/southasia}, were not
assigned to an alignment group, because a diaspora community, a general question
forum and two low-volume regional communities do not sit on the axis the
comparison runs along. Excluding them leaves the 305,217 comments across ten
communities in three alignment groups that \S\ref{sec:introduction} reports and
that every alignment-group result in this paper uses. Both figures describe the
same collection. Assigning communities is not assigning people, and no
community-level result here licenses a claim about an individual commenter.

\subsubsection{Collection and preprocessing}
\label{sec:data-collection}

% \begin{figure*}[t]
%   \centering
%   \includegraphics[width=\linewidth]{figures/s3_c1_collection_instrument}
%   \caption{The collection instrument, over the 49,494 collected submissions
%   (1 January 2019 to 10 May 2025). Panel A gives retrieval yield for the top
%   fourteen of thirty-nine queries. One query, \texttt{India Pakistan}, returned
%   21,655 submissions, or 44\% of the collection. Yields overlap, since 8,228
%   submissions (16.6\%) matched more than one query, so the bars do not sum to
%   the total. Panel B gives the three collection passes at the granularity each
%   was run at: a historical sweep at monthly resolution (36,767 submissions),
%   flashpoint windows at daily resolution (9,145), and the most recent crisis
%   window at hourly resolution (3,582). Each submission belongs to exactly one
%   pass. Collection stopped on 10 May 2025; the annotated corpus runs eight
%   months further, to 25 January 2026.}
%   \label{fig:collection}
% \end{figure*}

\candidate{S3.1.4}
The corpus holds 309,204 comments posted between January 2019 and January 2026,
a window spanning all six crisis events with margin on either side for the
temporal analyses. We retrieved them from the Arctic Shift
archive~\cite{arcticshift2026software} through the BAScraper
client~\cite{bascraper2026software}. Since the 2023 platform policy change
closed Pushshift to general use, working from a third-party archive is a
standing condition of this method rather than a footnote about
tooling~\cite{proferes2021studyingreddit}. Collection was query-based, and
Figure~\ref{fig:collection} makes the instrument itself visible rather than
asserted. Two properties of it bound what follows. The corpus is what a query
list found, and one query supplies close to half of it, so coverage is
concentrated in the vocabulary that query matches. And the three passes were run
at different temporal resolutions, so density in Panel B partly reflects
collection design rather than discourse volume alone. Panel B also marks a gap
we cannot close: collection stopped eight months before the corpus ends, so the
most recent months entered through a different path and the funnel rates below
do not describe them.

\candidate{S3.1.5}
Figure~\ref{fig:composition}B gives the six crisis windows, from Operation
Sindoor's 69,011 comments to the 2021 ceasefire's 3,923. Of the corpus, 161,446
comments fall inside those windows and 147,758 outside; no result in this paper
uses the outside portion. Because volume across events is uneven by more than an
order of magnitude, every event-level comparison we report is proportional and
made within a community, following the comparison rule established for
cross-crisis work~\cite{olteanu2015crisis}, and every headline result is
reported again with Operation Sindoor excluded. Preprocessing removed
duplicates, then any record whose body carried a removal or deletion marker, then
bot accounts, before a relevance gate that decided corpus membership and removed
59\% of what reached it. We did not remove records whose author field is
deleted, since deleting an account does not delete text already posted. Those
rates are reconstructed over the 219,056 comments with a raw record on disk, a
portion that excludes almost all of the two 2025 events. Dropping the marked
bodies, about a seventh of what survived deduplication, makes this a moderated
survivor set: what we analyse is what community moderation already let
stand~\cite{jhaver2019automoderator}. The reply analyses rebuild comment depth
from parent identifier chains rather than from the shipped column, and each
reply result names the denominator it used.